% abstract.tex
\chapter*{Abstract}

The transition from conceptual hardware sketching to software simulation represents a significant bottleneck in electronic design automation (EDA), traditionally demanding the tedious and manual drafting of SPICE netlists. This project presents an end-to-end, desktop-based machine learning pipeline engineered to automatically digitize hand-drawn circuit schematics into executable NGSPICE simulations. The core of the system leverages a YOLOv8 object detection architecture, fine-tuned via transfer learning to classify nine distinct electronic components while seamlessly handling extreme scale variance. To complement the deep learning model, a deterministic computer vision subsystem utilizes OpenCV for adaptive thresholding and morphological operations, algorithmically routing hand-drawn wire traces onto a discrete coordinate plane. These subsystems are integrated into a custom-built Graphical User Interface (GUI) developed in Python. This interface features an interactive digital canvas, allowing users to dynamically modify AI-detected components prior to the automated generation of the final \texttt{.cir} netlist. Quantitative evaluations, including precision-recall and confusion matrix analyses, demonstrate the model's high classification accuracy and minimal false-positive rate. Ultimately, this system effectively bridges the gap between rapid physical prototyping and rigorous digital simulation, substantially reducing schematic drafting overhead.
